\documentclass[11pt,a4paper,oneside]{report}
\usepackage[spanish,activeacute,es-noshorthands]{babel}
\usepackage[utf8]{inputenc}
\usepackage[T1]{fontenc}
\usepackage{lmodern}
\usepackage{amsmath, amssymb}
\usepackage{graphicx}
\usepackage[dvipsnames]{xcolor}
\usepackage{geometry}
\geometry{margin=2.4cm, headheight=14pt}
\usepackage{microtype}
\usepackage{booktabs}
\usepackage{longtable}
\usepackage{enumitem}
\usepackage{fancyhdr}
\usepackage{tcolorbox}
\tcbuselibrary{skins, breakable}
\usepackage{caption}
\captionsetup{font=small, labelfont=bf}
\usepackage{titlesec}
\usepackage{hyperref}
\hypersetup{
  colorlinks=true, linkcolor=NavyBlue, citecolor=ForestGreen, urlcolor=BrickRed,
  pdftitle={Resumen ejecutivo OreSat para INTISAT},
  pdfauthor={INTISAT}
}

\definecolor{intisat}{RGB}{26, 54, 93}
\definecolor{intisataccent}{RGB}{197, 143, 0}
\definecolor{ok}{RGB}{47, 133, 90}
\definecolor{warn}{RGB}{197, 60, 60}

\titleformat{\chapter}[hang]
  {\normalfont\Huge\bfseries\color{intisat}}
  {\thechapter.}{0.6em}{}
\titlespacing*{\chapter}{0pt}{-20pt}{20pt}

\pagestyle{fancy}
\fancyhf{}
\fancyhead[L]{\nouppercase{\leftmark}}
\fancyhead[R]{\thepage}
\renewcommand{\headrulewidth}{0.4pt}

\newtcolorbox{leccion}[1][]{
  colback=blue!4, colframe=intisat,
  fonttitle=\bfseries, title=Leccion aplicable al INTISAT,
  breakable, sharp corners=south, #1
}
\newtcolorbox{recomendacion}[1][]{
  colback=yellow!7, colframe=intisataccent,
  fonttitle=\bfseries, title=Recomendacion,
  breakable, sharp corners=south, #1
}
\newtcolorbox{caveat}[1][]{
  colback=red!4, colframe=warn,
  fonttitle=\bfseries, title=Cuidado / verificar,
  breakable, sharp corners=south, #1
}

\begin{document}

\begin{titlepage}
\centering
\vspace*{1cm}
{\color{intisat}\rule{\textwidth}{2pt}}
\vspace{0.6cm}

{\Huge\bfseries\color{intisat} OreSat}\\[0.4em]
{\LARGE\bfseries Resumen ejecutivo para el equipo INTISAT}

\vspace{0.7cm}
{\Large\itshape Que es, como esta construido, y que podemos\\
robar de su arquitectura}

\vspace{1cm}
{\color{intisataccent}\rule{0.6\textwidth}{1pt}}

\vspace{2cm}

\begin{tabular}{rl}
\textbf{Origen:} & Portland State Aerospace Society (PSAS) \\
\textbf{Universidad:} & Portland State University, Oregon, USA \\
\textbf{Filosofia:} & 100\% open-source (HW + FW + SW + GS) \\
\textbf{GitHub:} & \url{https://github.com/oresat} \\
\textbf{Audiencia:} & Lideres tecnicos del proyecto INTISAT \\
\textbf{Documento:} & v1.0 \\
\textbf{Fecha:} & \today
\end{tabular}

\vfill

\begin{minipage}{0.85\textwidth}
\textbf{Resumen.}\quad
OreSat es uno de los CubeSats academicos mas documentados y abiertos
del mundo. Construido por estudiantes de Portland State University, su
arquitectura modular por \emph{cards} comunicadas via CAN bus, su
documentacion publica, y su iteracion entre misiones (OreSat0, 0.5, 1)
lo hacen un caso de estudio ideal para el INTISAT. Este resumen identifica
las practicas mas valiosas que el equipo deberia adoptar, los repositorios
mas relevantes para revisar, y la comparacion directa contra el estado
actual del INTISAT.
\end{minipage}

\vspace{1.2cm}
{\color{intisat}\rule{\textwidth}{2pt}}
\end{titlepage}

\tableofcontents
\clearpage

% =====================================================================
\chapter{Que es OreSat}

\section{Origen y mision}

OreSat es el programa de CubeSats de \textbf{Portland State Aerospace
Society (PSAS)}, un grupo estudiantil de Portland State University
(Oregon, EEUU). El programa nacio en 2017 con un proposito doble:

\begin{itemize}[leftmargin=2em]
  \item Hacer ciencia espacial real (mediciones climaticas en LEO).
  \item Educar a estudiantes en ingenieria aeroespacial.
\end{itemize}

\section{La filosofia del open source}

Lo que diferencia a OreSat de otros CubeSats academicos:

\begin{itemize}[leftmargin=2em]
  \item Todo el \textbf{hardware} (esquematicos, layouts, BOMs) en GitHub.
  \item Todo el \textbf{firmware} (codigo de cada subsistema) en GitHub.
  \item Todo el \textbf{flight software} (OBC + apps) en GitHub.
  \item Todo el \textbf{ground station} (software, hardware, antenas) en GitHub.
  \item \textbf{Documentacion exhaustiva}: design reviews, lessons learned,
        test procedures.
\end{itemize}

Esto es raro. La mayoria de los CubeSats academicos liberan al menos
algunos componentes propietarios. OreSat libera \emph{todo}, en el
espiritu de la \emph{Libre Space Foundation} con la que tienen alianza
tecnica.

\section{Numeros del programa}

\begin{center}
\begin{tabular}{ll}
\toprule
\textbf{Atributo} & \textbf{Valor} \\
\midrule
Año de inicio                  & 2017 \\
Misiones lanzadas              & OreSat0 (2022) \\
Misiones en desarrollo         & OreSat0.5, OreSat1 \\
Tamaño del equipo              & 30-50 estudiantes activos \\
Tiempo de vuelo OreSat0        & 14 meses (deorbit enero 2024) \\
Repos en GitHub                & ~30 activos \\
Estandares aplicados           & CalPoly CDS, ISO 9001 \\
\bottomrule
\end{tabular}
\end{center}

% =====================================================================
\chapter{Las misiones}

\section{OreSat0 (2022): el primer vuelo}

Mision \textbf{1U} (un cubo de 10 cm) lanzada en mayo 2022 desde la ISS
via Nanoracks deployer. Objetivos:

\begin{itemize}[leftmargin=2em]
  \item Validar la arquitectura por cards en orbita real.
  \item Probar el OBC con Linux.
  \item Probar las comms UHF.
  \item Demostrar autonomia operativa.
\end{itemize}

\textbf{Resultado}: el satelite funciono nominalmente durante 14 meses
antes de quemarse en reentrada (enero 2024). Confirmo la viabilidad del
diseño.

\section{OreSat0.5: iteracion intermedia}

Version mejorada del OreSat0 con correcciones aprendidas. En desarrollo
al momento de redactar este documento. Mismo tamaño 1U.

\section{OreSat1: la mision cientifica}

CubeSat \textbf{2U} con payload cientifico principal: un imager de
\textbf{cirros (nubes altas)} para climatologia. Tambien lleva DXWiFi
(downlink de imagenes via WiFi modificado) y un livestream payload.

\begin{leccion}
OreSat siguio el patron \emph{primero validacion tecnica, despues mision
cientifica}. INTISAT podria adoptar el mismo: una mision corta validatoria,
seguida de la mision con payload de microscopia.
\end{leccion}

% =====================================================================
\chapter{Arquitectura modular --- las "cards"}

\section{El concepto de card}

OreSat se construye apilando \textbf{cards} (placas PCB) en formato
PC-104. Cada card es un subsistema independiente con:

\begin{itemize}[leftmargin=2em]
  \item Su propio microcontrolador.
  \item Su propio firmware.
  \item Su propio conector al backplane CAN.
\end{itemize}

\section{Las cards tipicas}

\begin{center}
\begin{tabular}{p{4cm}p{8cm}p{2cm}}
\toprule
\textbf{Card} & \textbf{Funcion} & \textbf{MCU} \\
\midrule
C3 (Command, Control, Comms) & OBC principal, Linux, comms UHF & ARM A9 \\
Solar                        & generacion solar + MPPT          & STM32 \\
Battery                      & gestion de bateria               & STM32 \\
Power                        & distribucion + reguladores       & STM32 \\
ADCS                         & control de actitud               & STM32 \\
Imager                       & camara payload                   & STM32 \\
DXWiFi                       & downlink por WiFi modificado     & STM32 \\
StarTracker                  & determinacion fina de actitud    & STM32 \\
\bottomrule
\end{tabular}
\end{center}

\section{La interconexion: CAN bus}

Todas las cards se comunican por un bus \textbf{CAN} compartido en el
backplane PC-104. Usan el protocolo \textbf{CANopen}, estandar industrial.

\begin{leccion}
Tu INTISAT usa I2C primario hacia el OBC. Si en el futuro tenes mas
cards (multiples sensores, subsystem propio para FPM), CAN+CANopen es
el camino que OreSat ya valido. Ahorras tiempo y errores tipicos.
\end{leccion}

\section{Por que esta arquitectura es buena}

\begin{itemize}[leftmargin=2em]
  \item \textbf{Cards intercambiables}: si falla la Battery card, se
        reemplaza sin afectar al resto.
  \item \textbf{Desarrollo paralelo}: cada equipo trabaja en su card
        independientemente.
  \item \textbf{Testing aislado}: cada card se testea sola en banco
        antes de integrar.
  \item \textbf{Reusabilidad entre misiones}: la Solar card de OreSat0
        se reuso en OreSat1.
  \item \textbf{Tolerancia a fallos}: una card en falla no contamina el
        bus CAN gracias a CANopen.
\end{itemize}

% =====================================================================
\chapter{Stack de software}

\section{Tres niveles}

OreSat tiene 3 capas de software:

\begin{itemize}[leftmargin=2em]
  \item \textbf{Firmware de cards}: bare-metal C o ChibiOS RTOS en STM32.
        Hace lo basico (sensar, actuar, hablar CAN).
  \item \textbf{Flight software del OBC (C3)}: Linux con servicios en
        Python. Coordina las cards, maneja comms, ejecuta logica de mision.
  \item \textbf{Ground station}: Python (oresat-gs), procesa telemetria y
        envia comandos.
\end{itemize}

\section{El C3 (Command, Control, Communications)}

Es el OBC principal de OreSat. Caracteristicas:

\begin{itemize}[leftmargin=2em]
  \item Linux Debian custom.
  \item Servicios Python: \texttt{oresat\_c3\_software}.
  \item Comunicacion con cards via socketCAN.
  \item Comms con tierra via UHF (modulo Si4463).
  \item Almacenamiento: eMMC.
\end{itemize}

\begin{leccion}
\textbf{Es exactamente lo que tu \texttt{cubesat/daemon.py} es para el
payload}, pero a escala de satelite completo. Vale leer el codigo del C3
para ver como organizan los servicios Python en un OBC operativo.
\end{leccion}

\section{Variantes con NASA cFS}

OreSat tambien explora la integracion con \textbf{NASA cFS} (Core Flight
System) en el repo \texttt{oresat0-cfs}. Esto es valioso para entender
como adaptar el framework profesional de NASA al hardware academico.

% =====================================================================
\chapter{Estructura de repositorios en GitHub}

Los 30+ repos de OreSat estan en \url{https://github.com/oresat}.
Los mas valiosos para revisar (por orden de prioridad):

\section{Prioridad alta para INTISAT}

\begin{itemize}[leftmargin=2em]
  \item \texttt{oresat-c3-software}: el flight software del OBC. \\
        \emph{Por que mirarlo}: tu daemon hace algo similar para el payload.
        Ver como organizan ellos un OBC entero.
  \item \texttt{oresat-cans}: definiciones del protocolo CAN. \\
        \emph{Por que}: si en el futuro agregas CAN como backup, esta es
        la plantilla.
  \item \texttt{oresat-firmware}: firmware de cada card. \\
        \emph{Por que}: ejemplos de como estructurar firmware embebido
        portable y testeable.
  \item \texttt{oresat-mainboard}: PCB del backplane. \\
        \emph{Por que}: como hacer un backplane PC-104 con CAN.
\end{itemize}

\section{Prioridad media}

\begin{itemize}[leftmargin=2em]
  \item \texttt{oresat-configs}: configuraciones de mision.
  \item \texttt{oresat-launch-services}: software ground station.
  \item \texttt{oresat0-imager}: payload de camara (similar concepto al
        del INTISAT pero EO en vez de microscopia).
  \item \texttt{oresat0-cfs}: integracion con NASA cFS.
\end{itemize}

\section{Prioridad baja (pero interesante)}

\begin{itemize}[leftmargin=2em]
  \item \texttt{oresat-pre-launch}: documentacion de las reviews.
  \item \texttt{oresat0-dxwifi}: downlink de imagenes via WiFi modificado
        (curioso pero no aplica directo).
  \item \texttt{oresat-livestream}: video downlink en orbita.
\end{itemize}

\begin{caveat}
La lista de repos puede haber cambiado desde la redaccion de este
documento. Verificar en \url{https://github.com/oresat} para el listado
actualizado.
\end{caveat}

% =====================================================================
\chapter{Lecciones aplicables al INTISAT}

Las practicas que mas vale la pena copiar:

\section{1. Documentacion como first-class deliverable}

OreSat publica TODO. Cada decision de diseño, cada lesson learned, cada
review tecnica. Eso les permite:

\begin{itemize}[leftmargin=2em]
  \item Onboardear miembros nuevos rapido.
  \item Recibir contribuciones externas.
  \item Justificar decisiones ante revisores.
  \item Reusar conocimiento entre misiones.
\end{itemize}

\textbf{Que adoptar}: tu \texttt{Documentos de Referencia/} esta bien
encaminado. Falta capturar las \emph{decisiones de diseño} (por que
elegimos X y no Y) y las \emph{lessons learned} (que aprendimos cuando
fallo Z).

\section{2. Arquitectura modular por cards}

Cada subsistema es independiente. Esto reduce la complejidad y permite
desarrollo paralelo.

\textbf{Que adoptar}: en el INTISAT, separar claramente los subsistemas
del payload. Tu payload ya esta bien aislado (carpeta \texttt{cubesat/}
con interfaz I2C clara). Mantener esa separacion estricta.

\section{3. CAN bus + CANopen para comunicacion entre cards}

OreSat eligio CAN porque es robusto, estandarizado, y soporta multi-master
naturalmente.

\textbf{Que adoptar}: cuando el INTISAT pase de OBC simple a multi-subsystem,
considerar CAN como bus primario (no I2C). I2C es mas simple pero menos
robusto a EMI.

\section{4. Iteracion entre misiones}

OreSat0 valida $\to$ OreSat0.5 mejora $\to$ OreSat1 hace ciencia. Cada
mision aprende de la anterior.

\textbf{Que adoptar}: planear el INTISAT como una serie, no como mision
unica. Primera mision: validar payload de microscopia con muestra simple.
Segunda mision: agregar funcionalidades (FPM completo, mas patrones de
OLED).

\section{5. Open source desde el dia 0}

OreSat NO empezo cerrado y luego abrio. Empezo abierto. Esto es una
disciplina dificil pero rinde:

\begin{itemize}[leftmargin=2em]
  \item Codigo mas limpio (saben que lo veran extraños).
  \item Pull requests externos.
  \item Visibilidad academica.
  \item Reproducibilidad.
\end{itemize}

\textbf{Que adoptar}: tu repo \texttt{CubeSat-EdgeAI-Payload} ya es publico.
Falta agregar:

\begin{itemize}[leftmargin=2em]
  \item Un \texttt{LICENSE} explicito (MIT o Apache 2.0).
  \item Un \texttt{CONTRIBUTING.md} con guidelines.
  \item Issues activos para que externos puedan contribuir.
  \item Codigo de conducta.
\end{itemize}

\section{6. Test-driven development para hardware embebido}

OreSat tiene tests automatizados de firmware (unit + integration). Esto
es raro en CubeSat academico pero crucial para confianza.

\textbf{Que adoptar}: tu daemon en Python deberia tener tests con pytest.
Especialmente para el dispatcher del I2C slave y para el OTA.

\section{7. Ground station propia y simple}

OreSat tiene su ground station en Portland. No depende de servicios de
terceros. Software Python facil de modificar.

\textbf{Que adoptar}: planear desde temprano una ground station academica
(antena UHF + SDR + Python con SatNOGS). No esperar al ultimo momento.

% =====================================================================
\chapter{Comparacion directa: OreSat vs INTISAT}

\begin{center}
\begin{tabular}{p{4cm}p{4.5cm}p{4.5cm}}
\toprule
\textbf{Aspecto} & \textbf{OreSat} & \textbf{INTISAT (actual)} \\
\midrule
Tamaño                   & 1U (OreSat0), 2U (OreSat1)     & 3U (planeado) \\
Filosofia                & 100\% open source               & open en GitHub, no formalizado \\
OBC                      & C3 card con Linux Debian        & Pi 5 (CM5 vuelo) con Pi OS \\
Buses internos           & CAN + CANopen                   & I2C (CAN como backup futuro) \\
Payload                  & Imager EO cirrus               & Microscopia FPM lensless \\
Cards                    & 7-8 cards independientes        & 1 card del payload + OBC central (separados) \\
Ground station           & propia, Python, en Portland     & no implementada aun \\
Misiones planeadas       & OreSat0/0.5/1, iterativo        & INTISAT (1 mision por ahora) \\
Documentacion publica    & exhaustiva, en GitHub           & creciendo, en \texttt{Documentos de Referencia/} \\
NASA cFS                 & explorado en \texttt{oresat0-cfs} & no usado (custom daemon) \\
Equipo                   & 30-50 estudiantes               & 4 personas \\
\bottomrule
\end{tabular}
\end{center}

\section{Donde estamos mejor que OreSat}

Pocas cosas, pero existen:

\begin{itemize}[leftmargin=2em]
  \item \textbf{Payload diferenciado}: microscopia biologica lensless es
        un nicho menos explotado que EO de cirros.
  \item \textbf{IA on-board}: usamos StarDist y Cellpose para procesar a
        bordo. OreSat solo captura.
  \item \textbf{Pipeline completo}: tenemos preprocesado, mejora con IA,
        segmentacion y exportacion en un solo daemon.
\end{itemize}

\section{Donde estamos atrasados}

\begin{itemize}[leftmargin=2em]
  \item \textbf{Tamaño del equipo}: ellos tienen 30-50, nosotros 4.
  \item \textbf{Misiones lanzadas}: ellos ya volaron 1, nosotros 0.
  \item \textbf{Modularidad}: ellos tienen 7 cards, nosotros 1.
  \item \textbf{Ground station}: ellos tienen, nosotros no.
  \item \textbf{Test coverage}: ellos automatizan, nosotros no.
\end{itemize}

% =====================================================================
\chapter{Plan de exploracion sugerido}

Para el equipo del INTISAT (4 personas), proponemos dividir la
investigacion de OreSat en 1 semana:

\begin{center}
\begin{tabular}{p{1.2cm}p{5.5cm}p{7cm}}
\toprule
\textbf{Dia} & \textbf{Repo / tema} & \textbf{Que extraer} \\
\midrule
1 & \texttt{oresat-c3-software} & arquitectura del OBC con Linux: como dividen servicios, como manejan estados \\
2 & \texttt{oresat-cans} + \texttt{oresat-firmware} & como definen CAN messages y los mapean a firmware \\
3 & \texttt{oresat0-imager} & comparar con tu payload de microscopia: que es comun, que es distinto \\
4 & \texttt{oresat-launch-services} & inspiracion para tu ground station \\
5 & Documentacion publica (wiki / docs) & lessons learned \\
\bottomrule
\end{tabular}
\end{center}

\section{Output esperado al final de la semana}

\begin{itemize}[leftmargin=2em]
  \item Lista de 5-10 practicas concretas que adoptar en INTISAT.
  \item Lista de 3-5 cosas a NO copiar (porque no aplican).
  \item Borrador de una arquitectura modular para INTISAT v2.
\end{itemize}

% =====================================================================
\chapter{Conclusion y proximos pasos}

OreSat es probablemente el \textbf{mejor caso de estudio publico}
para un CubeSat academico operativo. Vale la pena que el equipo de
INTISAT lo revise en profundidad.

\begin{recomendacion}
\textbf{Tres acciones concretas para esta semana}:

\begin{enumerate}[leftmargin=2em]
  \item Cada miembro del equipo clona un repo de OreSat segun la tabla
        del Capitulo 7 y lo explora 2 horas.
  \item Reunion del equipo (1 hr) para sintetizar: que adoptar, que no.
  \item Agregar \texttt{LICENSE} y \texttt{CONTRIBUTING.md} al
        repositorio del INTISAT (1 hora).
\end{enumerate}
\end{recomendacion}

\begin{caveat}
Este resumen ejecutivo se basa en informacion publica de OreSat al
momento de redaccion. Para detalles tecnicos especificos (nombres exactos
de repos, version de los protocolos, etc.) verificar siempre contra
\url{https://github.com/oresat} y \url{https://oresat.org}.
\end{caveat}

\section{Enlaces principales}

\begin{itemize}[leftmargin=2em]
  \item Organizacion GitHub: \url{https://github.com/oresat}
  \item Sitio web: \url{https://oresat.org}
  \item Portland State Aerospace Society: \url{https://psas.pdx.edu}
  \item Libre Space Foundation (aliados): \url{https://libre.space}
\end{itemize}

\end{document}
